% Part: incompleteness
% Chapter: computability-incompleteness
% Section: g1-halting-problem

\documentclass[../OpenLogic/include/open-logic-section]{subfiles}

\begin{document}

\olfileid{inc}{cpi}{hal}
\olsection{通过停机问题证明不完备性}

\begin{thm}\ollabel{thm:undecidable}
若 $\Th{T}$ 是一个 $\omega$-一致的理论，并且表示所有
原始递归关系，则 $\Th{T}$~是不可判定的。
\end{thm}

\begin{proof}
因为 $\Th{T}$ 表示所有原始递归函数，特别地，
它表示 Kleene 的 $T$ 谓词。也就是说，存在一个
公式~$!T(e, x, s)$，使得：
\begin{enumerate}
\item\ollabel{case:def} 若 $\cfind{e}(n) \fdefined$，则存在~$k \in \Nat$ 使得
$\Th{T} \Proves !T(\num{e}, \num{n}, \num{k})$，且
\item\ollabel{case:undef} 若 $\cfind{e}(n) \fundefined$，则对所有 $k \in \Nat$，
$\Th{T} \Proves \lnot !T(\num{e}, \num{n}, \num{k})$。
\end{enumerate}
我们将证明，若 $\Th{T}$ 是 $\omega$-一致的且表示
所有原始递归关系，我们就能判定
集合~$K = \Setabs{e}{\cfind{e}(e) \fdefined}$，即“自停机
问题”。

由于 $\Th{T}$ 是一个理论，它在后承下封闭。特别地，
若 $\Th{T} \Proves !T(\num{e}, \num{e}, \num{k})$，则
也有 $\Th{T} \Proves \lexists[z][!T(\num{e}, \num{e}, z)]$。结合
\olref{case:def}，这意味着若 $\cfind{e}(e) \fdefined$，
则 $\Th{T} \Proves \lexists[z][!T(\num{e}, \num{e}, z)]$。另一方面，
因为 $\Th{T}$ 是 $\omega$-一致的，我们不能同时对所有
$k \in \Nat$ 有 $\Th{T} \Proves \lnot !T(\num{e}, \num{e}, \num{k})$ 以及
$\Th{T} \Proves \lexists[z][!T(\num{e}, \num{e},
  z)]$。所以，由 \olref{case:undef}，若 $\cfind{e}(e) \fundefined$，则 $\Th{T} \Proves/
  \lexists[z][!T(\num{e}, \num{n}, z)]$。综上，
$\cfind{e}(e) \fdefined$ 当且仅当 $\Th{T} \Proves \lexists[z][!T(\num{e},
  \num{n}, z)]$。

如果 $\Th{T}$ 是可判定的，那么这就允许我们通过判定
是否 $\Th{T} \Proves \lexists[z][!T(\num{e}, \num{e}, z)]$ 来回答
$\cfind{e}(e)$ 是否有定义，即是否 $e \in K$。
然而，$K$ 是不可判定的（\olref[cmp][thy][ncp]{thm:K}）。
\end{proof}

\begin{thm}
  \ollabel{thm:complete-decidable}
  若\/ $\Th{T}$ 是一个完备、一致的可计算枚举理论，则
  $\Th{T}$ 是可判定的。
\end{thm}

\begin{proof}
  非形式地说，我们可以可计算地枚举一个可计算枚举理论的所有定理。
  要判定给定的 $!A$ 是否满足 $\Th{T} \Proves !A$，
  只需在这个定理枚举中搜索，直到找到 $!A$ 或~$\lnot !A$ 为止。
  由于 $\Th{T}$ 是完备的，这必定最终发生。
  如果在枚举中出现了 $!A$，我们就知道 $\Th{T} \Proves !A$。
  一致性保证了 $!A$ 和 $\lnot !A$ 中至多有一个是定理。
  于是，如果在枚举中出现了 $\lnot !A$，我们就知道 $\Th{T} \Proves/ !A$。

  更形式化地说，由于 $\Th{T}$ 是完备的、一致的，
  并且由一个可计算函数~$f$ 所枚举，
  $\Th{T}$ 的定理的哥德尔数集合的补集是
  \[\Complement{\Gn{\Th{T}}} = \Setabs{n}{\lexists[k][(\lnot
  \fn{Sent}(n) \lor (\Gn{\lnot} \concat f(k)) = n)]},\]（$\Gn{!A} \in
  \Nat \setminus \Gn{\Th{T}}$ 当且仅当 $!A$ 根本不是一个语句，
  或者其否定是~$\Th{T}$ 的定理）。
  根据 \olref[cmp][thy][eqc]{thm:exists-char}，
  $\Complement{\Gn{\Th{T}}}$ 是可计算枚举的，
  于是由 \olref[cmp][thy][cmp]{thm:ce-comp} 知
  $\Th{T}$ 是可判定的。
\end{proof}

\begin{cor}
  若\/ $\Th{T}$ 是 $\omega$-一致的、可计算枚举的，
  并且表示所有原始递归关系，则 $\Th{T}$ 是不完备的。
\end{cor}

\begin{proof}
  假设 $\Th{T}$ 是完备的。
  由于 $\omega$-一致的理论都是一致的，所以 $\Th{T}$ 是一致的，
  从而根据 \olref{thm:complete-decidable} 可知 $\Th{T}$ 是可判定的。
  这与 \olref{thm:undecidable} 矛盾。
\end{proof}

哥德尔第一不完全性定理的这个证明，并不能像通过对角化引理的证明那样
（参见 \olref[inc][inp][1in]{sec}）给出一个不可判定语句的具体例子。

\begin{thm}\ollabel{thm:g1}
  设 $\Th{T}$ 是一个 $\omega$-一致的、可计算枚举的理论，
  并且表示所有原始递归关系。
  则存在一个语句~$G$，使得 $\Th{T} \Proves/ G$ 且 $\Th{T} \Proves/ \lnot G$。
  $G$~可以从~$\Th{T}$ 的可计算枚举中明确指定。
\end{thm}

\begin{proof}
  设 $f$ 是~$\Th{T}$ 的一个可计算枚举。考虑函数~$g(e)$：若利用~$f$
  搜索 $\lnot\lexists[z][!T(\num{e}, \num{e}, z)]$ 的证明成功，则返回 $0$，否则无定义。令 $k$ 为~$g$ 的指标，
  即 $g = \cfind{k}$。那么，根据 $g$ 的定义，$g(e) \fdefined$ 当且仅当 $\Th{T} \Proves
  \lnot\lexists[z][!T(\num{e}, \num{e}, z)]$。
  因此存在 $m$ 使得 $T(k, e, m)$ 成立当且仅当 $\Th{T}
  \Proves \lnot\lexists[z][!T(\num{e}, \num{e}, z)]$。取 $e = k$，
  我们有：存在 $m$ 使得 $T(k, k, m)$ 成立当且仅当 $\Th{T}
  \Proves \lnot\lexists[z][!T(\num{k}, \num{k}, z)]$。

  假设 $g(k) \fdefined$，即 $\cfind{k}(k) \fdefined$，即对某个 $m$，$T(k, k, m)$ 成立。
  根据前述等价关系，我们有 $\Th{T} \Proves \lnot\lexists[z][!T(\num{k}, \num{k},
  z)]$。另一方面，由于 $\Th{T}$ 表示~$T$，
  我们有 $\Th{T} \Proves !T(\num{k}, \num{k}, \num{m})$，从而也有 $\Th{T}
  \Proves \lexists[z][!T(\num{k}, \num{k}, z)]$。这与~$\Th{T}$ 的一致性矛盾。
  我们可以断定 $g(k) \fundefined$，即 $\Th{T} \Proves/ \lnot\lexists[z][!T(\num{k}, \num{k}, z)]$。

  由于 $g(k) \fundefined$，对任何~$m$，$T(k, k, m)$ 都不成立。
  由于 $\Th{T}$ 表示~$T$，我们对所有~$m$ 都有 $\Th{T} \Proves \lnot
  !T(\num{k}, \num{k}, \num{m})$。
  因为 $\Th{T}$ 是 $\omega$-一致的，所以 $\Th{T} \Proves/ \lexists[z][!T(\num{k},
  \num{k}, z)]$。

  我们已经证明了，可以取语句 $G$ 为
  $\lnot\lexists[z][!T(\num{k}, \num{k}, z)]$。如果我们有~$g$ 的显式描述，就能确定~$g$ 的指标~$k$。
  再结合~$T$ 的显式描述以及~$T$ 在~$\Th{T}$ 中可由公式~$!T$ 表示的有效证明，
  我们就能从~$\Th{T}$ 的可计算枚举有效地找到 $G$。
\end{proof}

请记住，$\lexists[z][!T(\num{e}, \num{k}, z)]$ 说的是 $\cfind{e}(k)\fdefined$。因此，$\lexists[z][!T(\num{k},
\num{k}, z)]$ 说的是 $g(k) = \cfind{k}(k)\fdefined$，或者说
$g$ 自停机。换句话说，$g$ 在~$\Th{T}$ 中搜索这样一个公式的证明，该公式说的是
“$g$ 不自停机”；如果这样的证明存在，$g$ 就停机，否则不停机。

\end{document}